% =====================================================================
%  Two-frame Beamer deck for tv_pm/prior_routes/three_way_retrain.ipynb
%
%  Build:   pdflatex prior_routes_slides.tex
%
%  Figures: prior_recovery.pdf and threeway_cameraman.pdf, written by
%           prior_compare.py and compare_impl.py into tv_pm/results/figs/.
%           \graphicspath below finds them beside this file or at the repo path.
%
%  Numbers: the reported 250 000-step run (tv_pm/results/prior_recovery.csv and
%           threeway_psnr.csv). Reproduce with three_way_retrain.ipynb at
%           PRESET = "full"; its RUN_TAG/OUT_TAG keep a rerun off these files.
% =====================================================================
\documentclass[aspectratio=169,10pt]{beamer}

\usetheme{default}
\usecolortheme{seahorse}
\setbeamertemplate{navigation symbols}{}
\setbeamertemplate{footline}[frame number]
\setbeamerfont{frametitle}{size=\normalsize}
\setbeamersize{text margin left=4mm, text margin right=4mm}

\usepackage[T1]{fontenc}
\usepackage{graphicx}
\usepackage{booktabs}
\usepackage{amsmath,amssymb}
\usepackage{xcolor}

\graphicspath{{./}{./figs/}{../results/figs/}{../../results/figs/}}

\definecolor{iter}{HTML}{4C6EF5}
\definecolor{oneshot}{HTML}{F59F00}
\definecolor{direct}{HTML}{2F9E44}
\definecolor{bad}{HTML}{E03131}

\newcommand{\Iter}{\textcolor{iter}{\textbf{Iterative}}}
\newcommand{\Osh}{\textcolor{oneshot}{\textbf{One-shot}}}
\newcommand{\Dir}{\textcolor{direct}{\textbf{Direct fit}}}
\newcommand{\upm}{u_{\mathrm{PM}}}
\newcommand{\freg}{f_{\mathrm{reg}}}

\begin{document}

% =====================================================================
% FRAME 1 -- recovering the prior
% =====================================================================
\begin{frame}{Recovering a denoiser's implicit regularizer: three routes}
  \scriptsize

  $\upm$ is the prox of an implicit regularizer (Gribonval), so its optimality
  condition pins any correct prior on $\mathrm{range}(\upm)$ --- and no ground
  truth is needed:
  \vspace{-0.3em}
  \[
    \upm(x)=\mathrm{prox}_{\freg}(x)
    \;\iff\;
    \nabla \freg(y_k)=x_k-y_k,\quad y_k=\upm(x_k),
    \qquad\text{read against the noise floor }\delta .
  \]
  \vspace{-1.5em}

  \begin{columns}[T]
    \begin{column}{0.545\textwidth}
      \begin{tabular}{@{}p{0.20\linewidth}@{}p{0.78\linewidth}@{}}
        \toprule
        \Iter & $J_1=\langle x,w\rangle-\psi_\theta(w)-\tfrac12\|x\|^2$,\;
                $\nabla\psi_\theta(w)=x$ \;\textcolor{gray}{(inversion \emph{per query})}\\[0.2em]
        \Osh  & $J_2=G_\theta-\tfrac12\|x\|^2$,\; $G_\theta\approx\psi_\theta^{*}$
                \;\textcolor{gray}{(one forward pass)}\\[0.2em]
        \Dir  & $J_\theta$ fitted directly, $\nabla J_\theta(y_k)\approx x_k-y_k$
                \;\textcolor{gray}{(no recovery step)}\\
        \bottomrule
      \end{tabular}
      \vspace{0.25em}

      $J_1,J_2$ are \textbf{1-semiconvex for free} --- the $-\tfrac12\|\cdot\|^2$
      is in the Fenchel formula. $J_\theta$ is a plain ICNN, so it \emph{assumes}
      $\freg$ is convex: true for TV, false in general.
    \end{column}
    \begin{column}{0.445\textwidth}
      \centering
      \begin{tabular}{@{}lrrrr@{}}
        \toprule
        prior & resid & p90 & $\times\delta$ & cost \\
        \midrule
        \Iter & 14.48\% & \textcolor{bad}{371\%} & 6.6 & 12.7\,s \\
        \Osh  & 13.73\% & 98\% & 6.3 & 0.03\,s \\
        \Dir  & \textbf{8.07\%} & \textbf{21\%} & \textbf{3.7} & 0.03\,s \\
        \bottomrule
      \end{tabular}\\[0.2em]
      {\tiny 2000 held-out cameraman patches, $\delta=2.19\%$.\;
       $\nabla\psi_\theta$ vs $\upm$: \textbf{0.41\%} median --- everything above
       that is recovery-step error.}
    \end{column}
  \end{columns}

  \vspace{0.15em}
  \centering
  % top crop removes the figure's own suptitle (duplicated by the slide)
  \includegraphics[width=0.92\textwidth,trim=0 0 0 34,clip]{prior_recovery.pdf}

  \vspace{-0.2em}
  {\tiny \Iter{} has \emph{no safe} $\alpha$: at $\alpha=0$ the solve converges but
   0.45\% of patches leave the valid box; at $\alpha=0.1$ the $O(\alpha)$ bias takes
   it to 99.55\%.}

\end{frame}

% =====================================================================
% FRAME 2 -- the prior in action
% =====================================================================
\begin{frame}{The prior in action, and what the comparison establishes}
  \scriptsize

  \centering
  % top crop removes the figure's own 3-line suptitle
  \includegraphics[width=0.78\textwidth,trim=0 0 0 88,clip]{threeway_cameraman.pdf}

  \vspace{0.2em}
  \begin{columns}[T]
    \begin{column}{0.42\textwidth}
      \centering
      \begin{tabular}{@{}lrrr@{}}
        \toprule
        method & PSNR & vs $\upm$ & rel-$L^2$ \\
        \midrule
        noisy input & 22.50 & --- & --- \\
        $\upm$ \textcolor{gray}{(MCMC)} & \textbf{27.32} & --- & --- \\
        $u_\theta=\nabla\psi_\theta$ & 26.92 & 41.76 & 1.56\% \\
        $\mathrm{prox}_{J_2}$ \textcolor{gray}{(check)} & 26.83 & 39.50 & 2.02\% \\
        $\hat u=\mathrm{prox}_{J_\theta}$ & \textbf{27.09} & \textbf{46.93} & \textbf{0.86\%} \\
        \bottomrule
      \end{tabular}
    \end{column}
    \begin{column}{0.56\textwidth}
      \raggedright
      \textbf{Neither LPN route inverts anything to \emph{denoise}}: since
      $\nabla G=(\nabla\psi)^{-1}$, \Iter{} and \Osh{} both deploy
      $\nabla\psi_\theta$; they differ only in how they \emph{evaluate} the prior.
      The $\mathrm{prox}_{J_2}$ panel is a consistency check, not a denoiser --- for
      \Dir{} that object \emph{is} the deployed denoiser, so no gap can open.

      \vspace{0.35em}
      \textbf{1 --- the class is not negotiable} (mathematical). \Dir{} assumes
      $\freg$ convex; on a nonconvex prior it reads 101\% with curvature pinned at
      $0.000$ against an exact $-0.911$. No budget or dimension fixes it.

      \vspace{0.2em}
      \textbf{2 --- among valid routes, the ranking is a solver question}
      (numerical). Exact inversion (1-D): \Iter{} 0.02\% vs \Osh{} 1.22\%; Adam
      inversion at $d=64$: \Iter{} trails, at ${\sim}400\times$ the query cost.
    \end{column}
  \end{columns}

  \vspace{0.1em}
  {\tiny Cameraman, $\sigma=t=20/256$, $m=8000$, all nets 250\,000 steps.\;
   Reproduce: \texttt{three\_way\_retrain.ipynb}, \texttt{PRESET="full"}.}

\end{frame}

\end{document}
